%======================================================================
\Q{Problem 1: The four tests on a numeric two-port matrix}{\m{3+1+1+1+2}}
{\color{sub}\footnotesize \yr{\textbf{\texttt{79 Bh}}} \gap$\cdot$\gap the only fully numeric S-matrix in the archive}

\begin{asked}{\textbf{\texttt{79 Bh}} \ [3+1+1+1+2]}
Why are S-parameters used in high frequencies? The S-matrix of a certain microwave network is given as
$[S] = \begin{bmatrix} 0.4+j0.5 & j0.6 \\ j0.6 & 0.4-j0.5 \end{bmatrix}$.
(a) Is the network reciprocal? (b) Is the network lossless? (c) What is the return loss at the input? (d) If the input power to the network is 5 watts, what is the reflected power?
\end{asked}
\par\penalty0\Needspace{7\baselineskip}
\lead{Given}
\begin{itemize}
  \item $S_{11} = 0.4+j0.5$, \quad $S_{12} = S_{21} = j0.6$, \quad $S_{22} = 0.4-j0.5$.
  \item Input power $P_{in} = 5$ W at port 1; port 2 terminated in $Z_0$.
\end{itemize}
\lead{Step 0: the magnitudes, computed once and reused}
\begin{itemize}
  \item $|S_{11}| = \sqrt{0.4^2 + 0.5^2} = \sqrt{0.16+0.25} = \sqrt{0.41} = 0.6403$.
  \item $|S_{22}| = \sqrt{0.4^2 + (-0.5)^2} = \sqrt{0.41} = 0.6403$ (the same).
  \item $|S_{12}| = |S_{21}| = |j0.6| = 0.6$.
  \item Squares, which is what every power test needs:
    $|S_{11}|^2 = |S_{22}|^2 = 0.41$, \quad $|S_{12}|^2 = |S_{21}|^2 = 0.36$.
\end{itemize}
\par\penalty0\Needspace{14\baselineskip}
\lead{Step 1: run the four tests}
\begin{tabularx}{\linewidth}{@{}L{2.15cm}L{5.3cm}L{2.4cm}X@{}}
\toprule
\hrow \thead{Test} & \thead{Computation} & \thead{Result} & \thead{Conclusion} \\
\midrule
\textbf{Diagonal} & $S_{11} = 0.4+j0.5 \neq 0$, $S_{22} = 0.4-j0.5 \neq 0$
  & both non-zero & Neither port is matched \\
\textbf{Transpose} & $S_{12} = j0.6$ and $S_{21} = j0.6$
  & $S_{12} = S_{21}$ & \mk{(a) Reciprocal} \checkmark \\
\textbf{Column 1 sum} & $|S_{11}|^2 + |S_{21}|^2 = 0.41 + 0.36$
  & $0.77 \neq 1$ & \mk{(b) Not lossless} \\
\textbf{Column 2 sum} & $|S_{12}|^2 + |S_{22}|^2 = 0.36 + 0.41$
  & $0.77 \neq 1$ & Same shortfall at port 2 \\
\textbf{Column dot} & $S_{11}S_{12}^{*} + S_{21}S_{22}^{*}$
  $= (0.4+j0.5)(-j0.6) + (j0.6)(0.4+j0.5)$
  $= (0.30 - j0.24) + (-0.30 + j0.24)$ & $0$ & Phases are consistent; only the magnitudes fail \\
\bottomrule
\end{tabularx}
\vspace{2pt}
\begin{itemize}
  \item \textbf{(a)} $S_{12} = S_{21} = j0.6$, so the network is
    $\boxed{\text{reciprocal}}$: it is passive and contains no ferrite and no active
    device.
  \item \textbf{(b)} Both columns sum to $0.77$, not $1$, so the network is
    $\boxed{\text{not lossless}}$. It \textbf{dissipates} $1 - 0.77 = 0.23$, i.e.
    \textbf{23\%} of whatever is fed into either port.
  \item Worth adding: the column-dot test \emph{passes}. A network can have the right
    phase relations and still be lossy, which is why both halves of the unitary condition
    must be checked.
\end{itemize}
\lead{Step 2: (c) return loss at the input}
\begin{itemize}
  \item Return loss compares incident with reflected power at port 1, and $S_{11}$ is an
    \textbf{amplitude} ratio, so the factor is $20$:
    \[ RL = -20\log_{10}|S_{11}| = -20\log_{10}(0.6403) \]
  \item $\log_{10}(0.6403) = -0.19361$, so $RL = -20 \times (-0.19361)$.
  \item \textbf{(c)} $\boxed{RL = 3.87\ \text{dB}}$.
  \item Sanity check: a low return loss means a \textbf{bad} match, and indeed
    VSWR $= \dfrac{1+0.6403}{1-0.6403} = \dfrac{1.6403}{0.3597} = 4.56$, far from 1.
\end{itemize}
\lead{Step 3: (d) reflected power for 5 W in}
\begin{itemize}
  \item Power scales as the \textbf{square} of the wave amplitude:
    \[ P_{ref} = |S_{11}|^2 \, P_{in} = 0.41 \times 5 \]
  \item \textbf{(d)} $\boxed{P_{ref} = 2.05\ \text{W}}$.
\end{itemize}
\lead{Step 4: complete the power budget (free marks)}
\begin{tabularx}{\linewidth}{@{}L{3.4cm}L{4.4cm}L{2.0cm}X@{}}
\toprule
\hrow \thead{Component} & \thead{Formula} & \thead{Value} & \thead{Share of 5 W} \\
\midrule
Reflected at port 1 & $|S_{11}|^2 P_{in} = 0.41 \times 5$ & $2.05$ W & $41\%$ \\
Delivered to port 2 & $|S_{21}|^2 P_{in} = 0.36 \times 5$ & $1.80$ W & $36\%$ \\
Dissipated inside & $(1-0.77) \times 5$ & $1.15$ W & $23\%$ \\
\midrule
\textbf{Total} & & $\mathbf{5.00}$ W & $100\%$ \\
\bottomrule
\end{tabularx}
\vspace{2pt}
\begin{itemize}
  \item The three rows sum to the input, which is the arithmetic check that the whole
    answer is consistent.
  \item Insertion loss, if asked: $IL = -20\log_{10}(0.6) = \boxed{4.44\ \text{dB}}$.
\end{itemize}
\lead{Answer to the theory rider (3 marks)}
\begin{itemize}
  \item S-parameters are used at high frequencies because \textbf{total $V$ and $I$ are
    not uniquely defined} when the circuit is comparable to $\lambda$, because
    \textbf{open and short terminations cannot be realised} broadband and would
    \textbf{destroy or oscillate} active devices, and because \textbf{power} is what
    microwave instruments actually measure. See $\S$3.1.
\end{itemize}

\penalty0\vspace{2pt}

%======================================================================
\Q{Problem 2: Comparing two amplifier S-parameter sets}{\m{4}}
{\color{sub}\footnotesize \yr{\textbf{\texttt{82 Bh}}} \gap$\cdot$\gap $Z_0 = 50\,\Omega$; a reading question, not a design question}

\begin{asked}{\textbf{\texttt{82 Bh}} \ [4]}
Find the basic differences of the systems having the following sets of S-matrix.\\
(a) $S_{11} = 0.65\angle146^\circ$, $S_{12} = 0.12\angle46^\circ$, $S_{21} = 2.30\angle44^\circ$, $S_{22} = 0.17\angle{-172^\circ}$\\
(b) $S_{11} = 0.30\angle{-134^\circ}$, $S_{12} = 0.03\angle46^\circ$, $S_{21} = 2.20\angle44^\circ$, $S_{22} = 0.10\angle{-172^\circ}$
\end{asked}
\lead{Given}
\begin{itemize}
  \item Two \textbf{two-port} sets, quoted in polar form. Only the \textbf{magnitudes}
    matter for the comparison the question asks for; the angles matter only for
    $\Delta$ and $K$.
\end{itemize}
\lead{Step 1: what kind of device is each? (the two tests that decide it)}
\begin{tabularx}{\linewidth}{@{}L{2.15cm}L{4.4cm}L{3.3cm}X@{}}
\toprule
\hrow \thead{Test} & \thead{Computation} & \thead{Result} & \thead{Conclusion} \\
\midrule
\textbf{Transpose} & (a) $|S_{12}| = 0.12$ vs $|S_{21}| = 2.30$; (b) $0.03$ vs $2.20$
  & $S_{12} \neq S_{21}$ in both & \mk{Both are non-reciprocal} \\
\textbf{Column 1 sum} & (a) $0.65^2 + 2.30^2 = 0.4225 + 5.29$; (b) $0.30^2 + 2.20^2$
  & $5.71$ and $4.93$, both $\gg 1$ & \mk{Both are active}: they add power \\
\bottomrule
\end{tabularx}
\vspace{2pt}
\begin{itemize}
  \item $|S_{21}| > 1$ means the device \textbf{has gain}, so neither set can describe a
    passive component. Both are \textbf{transistor amplifiers} at the same frequency;
    the question is which is the better one.
\end{itemize}
\lead{Step 2: convert every magnitude into its engineering figure}
\begin{tabularx}{\linewidth}{@{}L{3.05cm}L{3.5cm}L{2.0cm}L{2.0cm}X@{}}
\toprule
\hrow \thead{Quantity} & \thead{Formula} & \thead{System (a)} & \thead{System (b)} & \thead{Better} \\
\midrule
Forward gain      & $20\log_{10}|S_{21}|$ & $7.23$ dB & $6.85$ dB & (a), barely \\
Reverse isolation & $-20\log_{10}|S_{12}|$ & $18.42$ dB & $30.46$ dB & \mk{(b), by 12 dB} \\
Input return loss & $-20\log_{10}|S_{11}|$ & $3.74$ dB & $10.46$ dB & \mk{(b)} \\
Input VSWR        & $\frac{1+|S_{11}|}{1-|S_{11}|}$ & $4.71$ & $1.86$ & \mk{(b)} \\
Output return loss & $-20\log_{10}|S_{22}|$ & $15.39$ dB & $20.00$ dB & (b) \\
Output VSWR       & $\frac{1+|S_{22}|}{1-|S_{22}|}$ & $1.41$ & $1.22$ & (b) \\
\bottomrule
\end{tabularx}
\vspace{2pt}
\begin{itemize}
  \item Sample line, so the arithmetic is visible:
    $-20\log_{10}(0.65) = -20 \times (-0.18709) = 3.74$ dB, and
    $\text{VSWR} = (1+0.65)/(1-0.65) = 1.65/0.35 = 4.71$.
\end{itemize}
\lead{Step 3: stability, the one place the angles are needed}
\begin{itemize}
  \item $\Delta = S_{11}S_{22} - S_{12}S_{21}$, and
    $K = \dfrac{1 - |S_{11}|^2 - |S_{22}|^2 + |\Delta|^2}{2|S_{12}S_{21}|}$.
  \item \textbf{System (a)}:
  \begin{itemize}
    \item $S_{11}S_{22} = (0.65)(0.17)\angle(146^\circ - 172^\circ)
      = 0.1105\angle{-26^\circ} = 0.0993 - j0.0484$.
    \item $S_{12}S_{21} = (0.12)(2.30)\angle(46^\circ + 44^\circ)
      = 0.276\angle90^\circ = j0.276$.
    \item $\Delta = 0.0993 - j0.0484 - j0.276 = 0.0993 - j0.3244$, so
      $|\Delta| = 0.3393$.
    \item $K = \dfrac{1 - 0.4225 - 0.0289 + 0.1151}{2 \times 0.276}
      = \dfrac{0.6637}{0.552} = 1.202$.
  \end{itemize}
  \item \textbf{System (b)}:
  \begin{itemize}
    \item $S_{11}S_{22} = (0.30)(0.10)\angle(-134^\circ - 172^\circ)
      = 0.03\angle54^\circ = 0.0176 + j0.0243$.
    \item $S_{12}S_{21} = (0.03)(2.20)\angle90^\circ = j0.066$.
    \item $\Delta = 0.0176 - j0.0417$, so $|\Delta| = 0.0453$.
    \item $K = \dfrac{1 - 0.09 - 0.01 + 0.0021}{2 \times 0.066}
      = \dfrac{0.9021}{0.132} = 6.834$.
  \end{itemize}
  \item Both have $K > 1$ and $|\Delta| < 1$, so both are \textbf{unconditionally
    stable}, but (b) is stable by a wide margin and (a) only just.
\end{itemize}
\lead{Answer: the basic differences}
\begin{itemize}
  \item Both are \textbf{active, non-reciprocal two-ports}, i.e. amplifiers of nearly the
    same gain ($7.23$ dB against $6.85$ dB, a difference of $0.4$ dB that is within normal
    device spread).
  \item \textbf{The real difference is everywhere except the gain}:
  \begin{itemize}
    \item \textbf{(b) is far better matched at the input}: $RL$ $10.46$ dB against
      $3.74$ dB, VSWR $1.86$ against $4.71$. System (a) reflects $42\%$ of the drive power
      straight back at the source.
    \item \textbf{(b) is far better isolated}: $|S_{12}| = 0.03$ against $0.12$, an extra
      $12$ dB. System (b) is close to \textbf{unilateral}, so it can be designed with the
      simple unilateral equations; system (a) needs the full bilateral treatment.
    \item \textbf{(b) is more stable}: $K = 6.83$ against $1.20$, and $|\Delta| = 0.045$
      against $0.339$. System (a) sits near the edge of unconditional stability, so a
      mismatched load could push it into oscillation.
  \end{itemize}
  \item \textbf{Verdict}: $\boxed{\text{(b) is the better amplifier}}$ at a cost of
    $0.4$ dB of gain: better match, better isolation, much better stability margin.
  \item \mk{Note for Chapter 6}: the same paper then asks for the maximum unilateral and
    bilateral gains of set (a). For the record they are $G_{TU,\max} = 9.75$ dB and
    $G_{\max} = 10.11$ dB; the derivation belongs to the amplifier chapter.
\end{itemize}

\penalty0\vspace{2pt}

%======================================================================
\Q{Problem 3: Evaluate the S-matrix of an amplitude modulator}{\m{7}}
{\color{sub}\footnotesize \yr{\textbf{78 Ch}} \gap$\cdot$\gap \added{} no source note carries this; modelled and evaluated here}

\begin{asked}{\textbf{78 Ch} \ [3+7]}
Elaborate the importance of S-parameters for microwave frequencies. Evaluate the S-matrix of an Amplitude Modulator.
\end{asked}
\lead{Step 1: model the modulator as a two-port}
\begin{itemize}
  \item A microwave amplitude modulator sits \textbf{in} the RF path: the carrier enters
    port 1 and the modulated carrier leaves port 2. The modulating signal is a
    \textbf{control input}, applied as bias, not as a third RF port.
  \item Physically it is a \textbf{PIN-diode variable attenuator}: the diode's RF
    resistance is set by its DC bias, so the modulating waveform steers the transmission
    through the device.
  \item Both RF ports are designed to stay \textbf{matched} across the modulation swing,
    otherwise the modulation would pull the source and produce incidental phase
    modulation.
\end{itemize}
\lead{Step 2: impose the three properties}
\begin{tabularx}{\linewidth}{@{}L{3.0cm}L{2.8cm}X@{}}
\toprule
\hrow \thead{Property} & \thead{Equation} & \thead{Why} \\
\midrule
Matched at both ports & $S_{11} = S_{22} = 0$
  & The attenuator is built as a matched pad at every bias point \\
Reciprocal & $S_{12} = S_{21} = \alpha$
  & Passive, no ferrite and no gain: a diode resistance is bilateral \\
Lossy, and \textbf{time varying} & $0 \le \alpha(t) \le 1$
  & The power removed from the carrier is dissipated in the diode \\
\bottomrule
\end{tabularx}
\vspace{2pt}
\begin{itemize}
  \item Hence
    \[ \boxed{\;[S] = \begin{bmatrix} 0 & \alpha(t) \\ \alpha(t) & 0 \end{bmatrix},
       \qquad \alpha(t) = \alpha_0\left(1 + m\cos\omega_m t\right)\;} \]
    with $\alpha_0$ the carrier-level transmission, $m$ the modulation index and
    $\omega_m$ the modulating frequency.
\end{itemize}
\lead{Step 3: check the condition}
\begin{itemize}
  \item Column sum: $|S_{11}|^2 + |S_{21}|^2 = \alpha^2$, which equals 1 only when
    $\alpha = 1$.
  \item So the modulator is \mk{never lossless} except at full transmission: the
    dissipated fraction is $1 - \alpha^2(t)$, and this is unavoidable. An amplitude
    modulator works by \textbf{throwing power away} in step with the modulating signal.
  \item It is reciprocal and matched, so the only non-ideal property is the loss.
\end{itemize}
\lead{Step 4: evaluate with numbers}
\begin{itemize}
  \item Take $\alpha_0 = 0.5$ and $m = 0.6$, a typical shallow modulation.
\end{itemize}
\begin{tabularx}{\linewidth}{@{}L{2.6cm}L{3.6cm}L{2.0cm}X@{}}
\toprule
\hrow \thead{Instant} & \thead{$\alpha(t)$} & \thead{$[S]$} & \thead{$IL$ and $P_{out}$ for $1$ W in} \\
\midrule
Peak, $\cos\omega_m t = +1$ & $0.5(1+0.6) = 0.80$
  & $\left[\begin{smallmatrix}0&0.8\\0.8&0\end{smallmatrix}\right]$
  & $IL = -20\log 0.8 = 1.94$ dB; $P_{out} = 0.64$ W \\
Carrier, $\cos\omega_m t = 0$ & $0.5$
  & $\left[\begin{smallmatrix}0&0.5\\0.5&0\end{smallmatrix}\right]$
  & $IL = 6.02$ dB; $P_{out} = 0.25$ W \\
Trough, $\cos\omega_m t = -1$ & $0.5(1-0.6) = 0.20$
  & $\left[\begin{smallmatrix}0&0.2\\0.2&0\end{smallmatrix}\right]$
  & $IL = 13.98$ dB; $P_{out} = 0.04$ W \\
\bottomrule
\end{tabularx}
\vspace{2pt}
\begin{itemize}
  \item \textbf{Recovering the modulation index from the matrix}, which is the check the
    examiner is looking for:
    \[ m = \frac{|S_{21}|_{\max} - |S_{21}|_{\min}}{|S_{21}|_{\max} + |S_{21}|_{\min}}
         = \frac{0.80 - 0.20}{0.80 + 0.20} = \frac{0.60}{1.00} = 0.6 \checkmark \]
  \item \textbf{Modulation depth in dB}: the transmission swings over
    $13.98 - 1.94 = 12.04$ dB.
\end{itemize}
\lead{Notes worth adding}
\begin{itemize}
  \item \textbf{Balanced realisation}: build the modulator on a magic tee, carrier into the
    H-arm, matched PIN diodes in the two collinear arms and the modulating bias on the
    E-arm. The two diodes reflect equally, the reflections recombine into the E-arm and
    the output amplitude follows the bias. The device is then a magic tee with two
    variable terminations, and the algebra is Ch4 companion Problem 12 with $\Gamma$ set by the bias
    instead of fixed at $-1$.
  \item \textbf{$S_{11}$ is the figure of merit}: a good modulator keeps $S_{11}$ small at
    \emph{every} bias point. If it does not, the source sees a varying load and the output
    picks up incidental phase modulation.
  \item \textbf{Contrast with an ideal switch}: $\alpha$ toggling between $1$ and $0$ makes
    the same matrix an on-off (ASK) modulator, with $m = 1$.
\end{itemize}

\penalty0\vspace{2pt}
